% ======================================================================
% GROUP 1: GENERAL (ข้อสอบจากไฟล์ Midterm เดิม)
% ======================================================================

\element{general}{
  \begin{question}{Q001}
    ข้อใดกล่าวถูกต้องเกี่ยวกับการเปลี่ยนแปลงทางดิจิทัล
    \begin{choices}
        \wrongchoice{Digitization หมายถึงการแปลงข้อมูลอนาล็อกเป็นรูปแบบดิจิทัลเท่านั้น}
        \wrongchoice{Digitalization หมายถึงการปรับปรุงและเปลี่ยนกระบวนการอนาล็อกให้เป็นดิจิทัล}
        \wrongchoice{Digital Transformation หมายถึงการปรับเปลี่ยนรูปแบบองค์กรโดยใช้เทคโนโลยีดิจิทัล}
        \correctchoice{ถูกทุกข้อ}
    \end{choices}
  \end{question}
}

\element{general}{
  \begin{question}{Q002}
    นิยามของเศรษฐศาสตร์คืออะไร
    \begin{choices}
        \wrongchoice{ศาสตร์ที่ศึกษาเฉพาะการผลิตสินค้าและบริการ}
        \wrongchoice{ศาสตร์ที่ศึกษาการผลิตและการกระจายสินค้าและบริการ}
        \correctchoice{ศาสตร์ที่ศึกษาการผลิต การกระจาย และการบริโภคสินค้าและบริการ}
        \wrongchoice{ผิดทุกข้อ}
    \end{choices}
  \end{question}
}

\element{general}{
  \begin{question}{Q003}
    ตามแนวคิดของ Joseph Schumpeter กราฟ S-Curve ใช้แสดงถึงสิ่งใด
    \begin{choices}
        \wrongchoice{วงจรชีวิตของสินค้าทั่วไป}
        \wrongchoice{เส้นต้นทุนของสินค้า}
        \wrongchoice{วัฏจักรความตื่นตัวของเทคโนโลยี}
        \correctchoice{วงจรชีวิตของนวัตกรรม}
    \end{choices}
  \end{question}
}

\element{general}{
  \begin{question}{Q004}
    ข้อใดเป็นตัวอย่างของ ``การทำลายล้างอย่างสร้างสรรค์'' (Creative Destruction)
    \begin{choices}
        \correctchoice{กล้องดิจิทัลและกล้องฟิล์ม}
        \wrongchoice{รถยนต์และเครื่องบิน}
        \wrongchoice{ส้อมกับมีดและตะเกียบ}
        \wrongchoice{ผิดทุกข้อ}
    \end{choices}
  \end{question}
}

\element{general}{
  \begin{question}{Q005}
    เหตุใดความก้าวหน้าของ S-curve จึงมักชะลอตัวในช่วงแรกก่อนเกิดการเติบโตอย่างก้าวกระโดด
    \begin{choices}
        \wrongchoice{ความอิ่มตัวของตลาด}
        \correctchoice{คอขวดทางเทคโนโลยี}
        \wrongchoice{อุปสรรคทางกฎระเบียบ}
        \wrongchoice{ข้อจำกัดด้านทรัพยากร}
    \end{choices}
  \end{question}
}

\element{general}{
  \begin{question}{Q006}
    เหตุใด S-curve จึงเติบโตอย่างรวดเร็วหลังจากผ่านจุด Breakthrough
    \begin{choices}
        \correctchoice{การค้นพบใหม่ ๆ นำไปสู่กระบวนการที่คล่องตัวและมีประสิทธิภาพมากขึ้น}
        \wrongchoice{นำไปสู่กฎระเบียบที่มากขึ้น}
        \wrongchoice{นำไปสู่กำไรที่มากขึ้น}
        \wrongchoice{ผิดทุกข้อ}
    \end{choices}
  \end{question}
}

\element{general}{
  \begin{question}{Q007}
    เหตุใดในที่สุด S-curve จึงต้องเผชิญกับผลตอบแทนลดน้อยถอยลง (Diminishing returns)
    \begin{choices}
        \wrongchoice{อุปสรรคทางกฎระเบียบ}
        \wrongchoice{การแข่งขันที่เพิ่มขึ้น}
        \correctchoice{ข้อจำกัดทางขีดความสามารถของเทคโนโลยี}
        \wrongchoice{ความต้องการของผู้บริโภคที่เปลี่ยนไป}
    \end{choices}
  \end{question}
}

\element{general}{
  \begin{question}{Q008}
    ทำไมกลุ่มโอเปก (OPEC) ถึงมีอำนาจมาก
    \begin{choices}
        \wrongchoice{มีสมาชิกหลายประเทศ}
        \wrongchoice{มีที่ดินเยอะ}
        \correctchoice{มีปริมาณน้ำมันดิบสำรองจำนวนมาก}
        \wrongchoice{ผิดทุกข้อ}
    \end{choices}
  \end{question}
}

\element{general}{
  \begin{question}{Q009}
    สิ่งใดเปรียบเสมือน "น้ำมันใหม่" (New Oil) แห่งยุคเศรษฐกิจดิจิทัล
    \begin{choices}
        \wrongchoice{บุคลากร (People)}
        \correctchoice{ข้อมูล (Data)}
        \wrongchoice{ความเร็วอินเทอร์เน็ต}
        \wrongchoice{เครือข่าย (Network)}
    \end{choices}
  \end{question}
}

\element{general}{
  \begin{question}{Q010}
    คำว่า FANGMAN หมายถึงอะไร
    \begin{choices}
        \correctchoice{กลุ่มบริษัทเทคโนโลยีชั้นนำที่มีรายได้และมูลค่าตลาดสูง}
        \wrongchoice{กลุ่มบริษัทพลังงานชั้นนำของโลก}
        \wrongchoice{ตัวย่อของกลุ่มประเทศมหาอำนาจ 7 ประเทศ}
        \wrongchoice{แวมไพร์}
    \end{choices}
  \end{question}
}

\element{general}{
  \begin{question}{Q011}
    ผลิตภัณฑ์มวลรวมในประเทศ (GDP) คืออะไร
    \begin{choices}
        \wrongchoice{มูลค่าตลาดของสินค้าและบริการ ``ขั้นกลาง'' ทั้งหมดในช่วงเวลาหนึ่ง}
        \correctchoice{มูลค่าตลาดของสินค้าและบริการ ``ขั้นสุดท้าย'' ทั้งหมดที่ผลิตในประเทศช่วงเวลาหนึ่ง}
        \wrongchoice{มาตรวัดอรรถประโยชน์ของสินค้าขั้นกลางทั้งหมด}
        \wrongchoice{มาตรวัดอรรถประโยชน์ของสินค้าขั้นสุดท้ายทั้งหมด}
    \end{choices}
  \end{question}
}

\element{general}{
  \begin{question}{Q012}
    นิยามของ ``สินค้าขั้นสุดท้าย'' (Final Goods) คืออะไร
    \begin{choices}
        \wrongchoice{สินค้าที่ผลิตเสร็จเป็นลำดับสุดท้ายของสายการผลิต}
        \wrongchoice{สินค้าที่ใกล้หมดอายุ}
        \wrongchoice{สินค้าที่ถูกนำไปใช้เป็นวัตถุดิบเพื่อผลิตต่อ}
        \correctchoice{สินค้าที่ผู้บริโภคซื้อไปเพื่อการบริโภคหรือใช้งานโดยตรง}
    \end{choices}
  \end{question}
}

\element{general}{
  \begin{question}{Q013}
    ข้อแตกต่างระหว่าง GDP ที่เป็นตัวเงิน (Nominal) และ GDP ที่แท้จริง (Real) คืออะไร
    \begin{choices}
        \wrongchoice{Nominal GDP มีความแม่นยำกว่า Real GDP}
        \correctchoice{Real GDP คือ Nominal GDP ที่ปรับลดผลกระทบจากเงินเฟ้อออกไปแล้ว}
        \wrongchoice{Real GDP คือ Nominal GDP ที่ปรับลดผลกระทบจากจำนวนประชากรแล้ว}
        \wrongchoice{ผิดทุกข้อ}
    \end{choices}
  \end{question}
}

\element{general}{
  \begin{question}{Q014}
    ตามแบบจำลอง Solow ข้อใด \textbf{ไม่ใช่} ปัจจัยหลักที่ส่งผลต่อ Real GDP
    \begin{choices}
        \wrongchoice{ทุน (Capital)}
        \wrongchoice{แรงงานที่มีประสิทธิภาพ (Effective Labor)}
        \correctchoice{ผลิตภัณฑ์ (Product)}
        \wrongchoice{นวัตกรรม (Innovation)}
    \end{choices}
  \end{question}
}

\element{general}{
  \begin{question}{Q015}
    หากเรามีเงินทุน (Capital) มากกว่าระดับ Steady-state เราควรทำอย่างไร
    \begin{choices}
        \wrongchoice{เพิ่มทุนเพราะต้องลดค่าเสื่อมราคา}
        \wrongchoice{เพิ่มทุนเพราะจะสร้างผลตอบแทนได้มากขึ้น}
        \correctchoice{ลดทุนเพราะค่าเสื่อมราคาจะสูงกว่าการลงทุนใหม่}
        \wrongchoice{ลดทุนเพราะจะสร้างผลตอบแทนได้มากขึ้น}
    \end{choices}
  \end{question}
}

\element{general}{
  \begin{question}{Q016}
    ดัชนีจีนี (Gini Index) ใช้สำหรับวัดสิ่งใด
    \begin{choices}
        \wrongchoice{ความสุขของประชากร}
        \wrongchoice{อายุขัยเฉลี่ย}
        \wrongchoice{ระดับการศึกษา}
        \correctchoice{ความเหลื่อมล้ำในการกระจายรายได้หรือความมั่งคั่ง}
    \end{choices}
  \end{question}
}

% \element{general}{
%   \begin{question}{Q017}
%     ตามทฤษฎีปริมาณเงิน (QTM) หากรัฐบาลกระตุ้นให้ Velocity เพิ่มขึ้น โดยปริมาณเงินเท่าเดิม จะเกิดผลอย่างไร
%     \begin{choices}
%         \correctchoice{อัตราเงินเฟ้อเพิ่มขึ้น}
%         \wrongchoice{GDP ที่แท้จริงลดลง}
%         \wrongchoice{อัตราการว่างงานเพิ่มขึ้น}
%         \wrongchoice{GDP ตัวเงินลดลง}
%     \end{choices}
%   \end{question}
% }

% \element{general}{
%   \begin{question}{Q018}
%     ข้อใด \textbf{ไม่ใช่} สาเหตุของเงินเฟ้อ
%     \begin{choices}
%         \wrongchoice{Cost push}
%         \wrongchoice{Demand Pull}
%         \wrongchoice{Built-in}
%         \correctchoice{Unemployment (การว่างงาน)}
%     \end{choices}
%   \end{question}
% }

\element{general}{
  \begin{question}{Q019}
    หน่วยย่อยที่สุดของข้อมูลดิจิทัลในคอมพิวเตอร์คืออะไร
    \begin{choices}
        \wrongchoice{Software}
        \wrongchoice{Hardware}
        \wrongchoice{Byte}
        \correctchoice{Bit}
    \end{choices}
  \end{question}
}

\element{general}{
  \begin{question}{Q020}
    ข้อแตกต่างหลักระหว่าง Digital และ Analog คืออะไร
    \begin{choices}
        \wrongchoice{Analog คือเก่า Digital คือใหม่}
        \correctchoice{Analog มีความต่อเนื่อง ส่วน Digital มีความไม่ต่อเนื่อง}
        \wrongchoice{วัสดุที่ใช้ทำอุปกรณ์}
        \wrongchoice{ความชอบของผู้ใช้งาน}
    \end{choices}
  \end{question}
}

\element{general}{
  \begin{question}{Q021}
    จุดประสงค์ของทรานซิสเตอร์ (Transistors) ในคอมพิวเตอร์คืออะไร
    \begin{choices}
        \correctchoice{เก็บข้อมูลบิต}
        \wrongchoice{แปลงข้อมูลดิจิทัล}
        \wrongchoice{แปลงไฟฟ้า AC เป็น DC}
        \wrongchoice{แปลงไฟฟ้า DC เป็น AC}
    \end{choices}
  \end{question}
}

\element{general}{
  \begin{question}{Q022}
    ข้อใดคืออุปกรณ์ Output ที่แปลงข้อมูลดิจิทัลกลับมาเป็นข้อมูลที่มนุษย์รับรู้ได้
    \begin{choices}
        \wrongchoice{ไมโครโฟน}
        \wrongchoice{คีย์บอร์ด}
        \correctchoice{หน้าจอแสดงผล}
        \wrongchoice{ทัชแพด}
    \end{choices}
  \end{question}
}

\element{general}{
  \begin{question}{Q023}
    จุดประสงค์ของตาราง ASCII คืออะไร
    \begin{choices}
        \correctchoice{แปลงสัญลักษณ์ตัวอักษรเป็นตัวเลข}
        \wrongchoice{แปลงสัญลักษณ์ตัวอักษรเป็นรหัสมอร์ส}
        \wrongchoice{แปลงสัญญาณอนาล็อกเป็นตัวเลข}
        \wrongchoice{แปลงสัญญาณอนาล็อกเป็นรหัสมอร์ส}
    \end{choices}
  \end{question}
}

\element{general}{
  \begin{question}{Q024}
    ความแตกต่างระหว่างโมเดลสี RGB และ CMYK คืออะไร
    \begin{choices}
        \correctchoice{RGB เป็นโมเดลสีของแสง ส่วน CMYK เป็นโมเดลสีของหมึกพิมพ์}
        \wrongchoice{จำนวนพิกเซลที่เก็บได้}
        \wrongchoice{จำนวนช่องสี}
        \wrongchoice{ไม่มีข้อแตกต่าง}
    \end{choices}
  \end{question}
}

% \element{general}{
%   \begin{question}{Q025}
%     เหตุใดการแปลงรสและกลิ่นเป็นดิจิทัลจึงยากกว่าภาพและเสียง
%     \begin{choices}
%         \wrongchoice{เพราะคุณสมบัติทางกายภาพมีมากเกินไป}
%         \correctchoice{เพราะเกิดจากปฏิกิริยาทางเคมี}
%         \wrongchoice{เพราะรสชาติไม่มีหน่วยย่อย}
%         \wrongchoice{ถูกทุกข้อ}
%     \end{choices}
%   \end{question}
% }

\element{general}{
  \begin{question}{Q026}
    Miyashita Lan (Meiji University) ประสบความสำเร็จในการแปลงรสชาติใดเป็นดิจิทัล
    \begin{choices}
        \wrongchoice{หวาน}
        \wrongchoice{เปรี้ยว}
        \correctchoice{เค็ม}
        \wrongchoice{อูมามิ}
    \end{choices}
  \end{question}
}

% \element{general}{
%   \begin{question}{Q027}
%     จุดประสงค์ของ Dataland (MIT 1970s) คืออะไร
%     \begin{choices}
%         \wrongchoice{สวนสนุกคอมพิวเตอร์}
%         \correctchoice{ให้คนโต้ตอบกับคอมพิวเตอร์ผ่านระบบจัดการข้อมูลเชิงพื้นที่ (Spatial data-management)}
%         \wrongchoice{ใช้เก็บข้อมูลสำคัญ}
%         \wrongchoice{ผิดทุกข้อ}
%     \end{choices}
%   \end{question}
% }

% \element{general}{
%   \begin{question}{Q028}
%     จุดประสงค์ของ Neuralink คืออะไร
%     \begin{choices}
%         \wrongchoice{ควบคุมคอมพิวเตอร์ด้วยสัญญาณสมอง}
%         \wrongchoice{สื่อสารผ่านกระแสจิต}
%         \wrongchoice{พัฒนาการฝังชิปในสมอง}
%         \correctchoice{ถูกทุกข้อ}
%     \end{choices}
%   \end{question}
% }

% \element{general}{
%   \begin{question}{Q029}
%     สัตว์ชนิดใดที่ถูกใช้ทดสอบ Neuralink เล่นเกม Pong
%     \begin{choices}
%         \wrongchoice{หมู}
%         \wrongchoice{หนู}
%         \correctchoice{ลิง}
%         \wrongchoice{ไก่}
%     \end{choices}
%   \end{question}
% }

\element{general}{
  \begin{question}{Q030}
    ความแตกต่างระหว่าง Web 1.0 และ Web 2.0 คืออะไร
    \begin{choices}
        \correctchoice{Web 1.0 เป็น Static (อ่านอย่างเดียว), Web 2.0 เป็น Dynamic (โต้ตอบได้)}
        \wrongchoice{Web 1.0 เป็น Dynamic, Web 2.0 เป็น Static}
        \wrongchoice{เหมือนกัน}
        \wrongchoice{ผิดทุกข้อ}
    \end{choices}
  \end{question}
}

\element{general}{
  \begin{question}{Q031}
    ข้อใดคือตัวอย่างของแพลตฟอร์มตัวกลาง (Intermediary Platform)
    \begin{choices}
        \wrongchoice{ระบบภายในบริษัท}
        \wrongchoice{เว็บสายการบินโดยตรง}
        \wrongchoice{บล็อกส่วนตัว}
        \correctchoice{Traveloka}
    \end{choices}
  \end{question}
}

\element{general}{
  \begin{question}{Q032}
    โมเดลธุรกิจใดทำได้ยากที่สุดหากไม่มีแพลตฟอร์มดิจิทัล
    \begin{choices}
        \wrongchoice{B2C}
        \wrongchoice{B2B}
        \correctchoice{C2C}
        \wrongchoice{B2G}
    \end{choices}
  \end{question}
}

% \element{general}{
%   \begin{question}{Q033}
%     CBEC คืออะไร
%     \begin{choices}
%         \correctchoice{การค้าขายออนไลน์ข้ามพรมแดนถึงผู้บริโภคโดยตรง}
%         \wrongchoice{อีคอมเมิร์ซในประเทศ}
%         \wrongchoice{ระบบรัฐต่อรัฐ}
%         \wrongchoice{ตลาดนัดท้องถิ่น}
%     \end{choices}
%   \end{question}
% }

\element{general}{
  \begin{question}{Q034}
    ข้อใดเป็นต้นทุนคงที่ (Fixed Cost)
    \begin{choices}
        \wrongchoice{ค่าคอมมิชชั่น}
        \wrongchoice{ค่าไฟการผลิต}
        \correctchoice{เงินเดือน CEO}
        \wrongchoice{ค่าวัตถุดิบ}
    \end{choices}
  \end{question}
}

\element{general}{
  \begin{question}{Q035}
    แพลตฟอร์มช่วยลดต้นทุนคงที่ได้อย่างไร
    \begin{choices}
        \wrongchoice{ลดต้นทุนธุรกรรม}
        \correctchoice{ลดความจำเป็นในการมีหน้าร้านจริง}
        \wrongchoice{ลดค่าวัตถุดิบ}
        \wrongchoice{ถูกทุกข้อ}
    \end{choices}
  \end{question}
}

\element{general}{
  \begin{question}{Q036}
    จากตารางต้นทุน (Table A) ค่าของ \textbf{X} (ATC) คือเท่าใด

    \begin{center}
      \small
      \setlength{\tabcolsep}{4pt}
      \begin{tabular}{c|c|c|c|c}
        \textbf{L} & \textbf{Q} & \textbf{TC} & \textbf{ATC} & \textbf{MC} \\ \hline
        2 & 10 & 4,400 & \textbf{X} & - \\
        3 & 12 & 4,700 & 391.6 & \textbf{Y}
      \end{tabular}
      
      \vspace{0.1cm}
      \scriptsize
      (L=แรงงาน, Q=ผลผลิต, TC=ต้นทุนรวม, ATC=ต้นทุนเฉลี่ย, MC=ต้นทุนส่วนเพิ่ม)
    \end{center}

    \begin{choices}
        \wrongchoice{435}
        \correctchoice{440}
        \wrongchoice{460}
        \wrongchoice{500}
    \end{choices}
  \end{question}
}

\element{general}{
  \begin{question}{Q037}
    พิจารณาตารางต้นทุนการผลิตด้านล่าง \\
    ค่าของ \textbf{Y} (ต้นทุนเพิ่ม MC เมื่อเพิ่มแรงงานเป็น 3 คน) คือเท่าใด

    \begin{center}
      \small
      \setlength{\tabcolsep}{4pt}
      \begin{tabular}{c|c|c|c}
        \textbf{แรงงาน (L)} & \textbf{ผลผลิต (Q)} & \textbf{TC} & \textbf{MC} \\ \hline
        2 & 20 & 2,000 & - \\
        3 & 22 & 2,300 & \textbf{Y}
      \end{tabular}
    \end{center}

    \begin{choices}
        \wrongchoice{300}    % ตัวลวง: เอาส่วนต่าง TC มาตอบเลย (ลืมหาร Q)
        \wrongchoice{250}
        \wrongchoice{176}
        \correctchoice{150}  % คำตอบถูก: (2300-2000) / (22-20) = 300/2
    \end{choices}
  \end{question}
}

\element{general}{
  \begin{question}{Q038}
    ผู้ผลิตจะได้กำไรสูงสุดเมื่อผลิตที่จุดใด (พิจารณา Marginal Cost)
    \begin{choices}
        \correctchoice{MC = Price}
        \wrongchoice{MC > Price}
        \wrongchoice{MC < Price}
        \wrongchoice{ได้กำไรเสมอ}
    \end{choices}
  \end{question}
}

\element{general}{
  \begin{question}{Q039}
    ตัวอย่างตลาดที่ใกล้เคียงกับตลาดแข่งขันสมบูรณ์ (Perfect Competition) มากที่สุด
    \begin{choices}
        \wrongchoice{นาฬิกาหรู}
        \wrongchoice{คอนโดมิเนียม}
        \wrongchoice{เสื้อผ้า}
        \correctchoice{ข้าวคุณภาพต่ำ}
    \end{choices}
  \end{question}
}

\element{general}{
  \begin{question}{Q040}
    ความเสี่ยงหลักของการใช้แพลตฟอร์มของบริษัทอื่นคืออะไร
    \begin{choices}
        \wrongchoice{ราคาขึ้น}
        \wrongchoice{เสียการควบคุมลูกค้า}
        \wrongchoice{แพลตฟอร์มดึงคู่แข่งมาแข่ง}
        \correctchoice{ถูกทุกข้อ}
    \end{choices}
  \end{question}
}

% \element{general}{
%   \begin{question}{Q041}
%     สินค้าที่มีสภาพเป็นปรปักษ์ (Rivalrous) และกีดกันได้ (Excludable) คือสินค้าประเภทใด
%     \begin{choices}
%         \correctchoice{Private Goods}
%         \wrongchoice{Common Goods}
%         \wrongchoice{Club Goods}
%         \wrongchoice{Public Goods}
%     \end{choices}
%   \end{question}
% }

\element{general}{
  \begin{question}{Q042}
    ข้อได้เปรียบหลักของสินค้าสารสนเทศ (Information Goods) คืออะไร
    \begin{choices}
        \wrongchoice{การตลาด}
        \wrongchoice{การละเมิดลิขสิทธิ์}
        \correctchoice{ความสามารถในการขยายขนาด}
        \wrongchoice{ความปลอดภัย}
    \end{choices}
  \end{question}
}

\element{general}{
  \begin{question}{Q043}
    ข้อใด \textbf{ไม่ใช่} ประโยชน์ของการใช้ดิจิทัลแพลตฟอร์ม (ในบริบทลดต้นทุน)
    \begin{choices}
        \wrongchoice{ลดต้นทุนคงที่}
        \wrongchoice{ลดต้นทุนธุรกรรม}
        \wrongchoice{ลดต้นทุนขนส่ง}
        \correctchoice{การแบ่งปันความรู้}
    \end{choices}
  \end{question}
}

% \element{general}{
%   \begin{question}{Q044}
%     ข้อแตกต่างสำคัญระหว่าง Metaverse กับเกมออนไลน์ปัจจุบันคืออะไร
%     \begin{choices}
%         \wrongchoice{ความเป็นเจ้าของไอเทม}
%         \wrongchoice{VR}
%         \correctchoice{อวตารและไอเทมสามารถย้ายข้ามแพลตฟอร์มได้ (Interoperability)}
%         \wrongchoice{Mass multiplayer}
%     \end{choices}
%   \end{question}
% }

\element{general}{
  \begin{question}{Q045}
    Hype Cycle คืออะไร
    \begin{choices}
        \correctchoice{วงจรชีวิตของนวัตกรรม}
        \wrongchoice{วงจรเฉพาะที่เป็นไวรัล}
        \lastchoices
        \wrongchoice{ถูกทั้งสองข้อแรก}
        \wrongchoice{ผิดทั้งสองข้อแรก}
    \end{choices}
  \end{question}
}

% \element{general}{
%   \begin{question}{Q046}
%     ทำไม Metaverse ถึงต้องการ Web 3.0
%     \begin{choices}
%         \wrongchoice{เร็วกว่า}
%         \wrongchoice{เหมาะกับ VR}
%         \wrongchoice{รองรับคนเยอะ}
%         \correctchoice{เพื่อความเป็นเจ้าของคอนเทนต์แบบกระจายศูนย์ (Decentralized ownership)}
%     \end{choices}
%   \end{question}
% }

% \element{general}{
%   \begin{question}{Q047}
%     Digital Twin คืออะไร
%     \begin{choices}
%         \correctchoice{แบบจำลองเสมือนของสินทรัพย์ทางกายภาพที่แสดงข้อมูลเรียลไทม์}
%         \wrongchoice{อวตารดิจิทัล}
%         \wrongchoice{เงินดิจิทัล}
%         \wrongchoice{ไม่มีข้อถูก}
%     \end{choices}
%   \end{question}
% }

% \element{general}{
%   \begin{question}{Q048}
%     Interoperability ด้านใดที่จำเป็นสำหรับการถ่ายโอนความเป็นเจ้าของ (Ownership)
%     \begin{choices}
%         \wrongchoice{Behavior}
%         \wrongchoice{Meaning}
%         \wrongchoice{Presentation}
%         \correctchoice{Persistence}
%     \end{choices}
%   \end{question}
% }

% \element{general}{
%   \begin{question}{Q049}
%     Interoperability ด้านใดที่เกี่ยวกับรูปแบบกราฟิก (Graphic style)
%     \begin{choices}
%         \wrongchoice{Behavior}
%         \wrongchoice{Meaning}
%         \correctchoice{Presentation}
%         \wrongchoice{Persistence}
%     \end{choices}
%   \end{question}
% }

% \element{general}{
%   \begin{question}{Q050}
%     Interoperability ด้านใดที่เกี่ยวกับการตีความคุณสมบัติไอเทม (Interpret property)
%     \begin{choices}
%         \wrongchoice{Behavior}
%         \correctchoice{Meaning}
%         \wrongchoice{Presentation}
%         \wrongchoice{Persistence}
%     \end{choices}
%   \end{question}
% }

% \element{general}{
%   \begin{question}{Q051}
%     หากแบ่ง Metaverse แบบเศรษฐกิจปิด+ภายใน (Closed + Internalized) จะเรียกว่าอะไร
%     \begin{choices}
%         \wrongchoice{Sandboxes}
%         \wrongchoice{Walled Gardens}
%         \correctchoice{Theme Parks}
%         \wrongchoice{Mashup Ecosystems}
%     \end{choices}
%   \end{question}
% }

% \element{general}{
%   \begin{question}{Q052}
%     โมเดลธุรกิจใดมีความเป็นไปได้ต่ำที่สุดสำหรับ Metaverse
%     \begin{choices}
%         \correctchoice{Premium game (ซื้อขาด)}
%         \wrongchoice{Free-to-play}
%         \wrongchoice{Play-to-earn}
%         \wrongchoice{Advertising}
%     \end{choices}
%   \end{question}
% }

\element{general}{
  \begin{question}{Q053}
    ลักษณะสำคัญของระบบเศรษฐกิจแบบ \textbf{Gift Economy} คืออะไร
    \begin{choices}
        \correctchoice{การให้โดยไม่หวังผลตอบแทนในทันที (เน้นสายสัมพันธ์)}
        \wrongchoice{การใช้บัตรของขวัญ (Gift Voucher) แทนเงินสด}
        \wrongchoice{การแลกเปลี่ยนสินค้ากันโดยตรง (Barter System)}
        \wrongchoice{การซื้อขายสินค้าออนไลน์}
    \end{choices}
  \end{question}
}

\element{general}{
  \begin{question}{Q054}
    ในยุคโบราณที่ยังไม่มีเงินตรา ชาวบ้านนำไก่ไปแลกกับข้าวสาร \\
    ระบบเศรษฐกิจแบบนี้เรียกว่า \textbf{Barter Economy} ซึ่งหมายถึงอะไร
    \begin{choices}
        \wrongchoice{เศรษฐกิจการต่อรองราคา}
        \wrongchoice{ระบบที่ใช้เงินเป็นสื่อกลาง}
        \correctchoice{การแลกเปลี่ยนสินค้ากับสินค้า}
        \wrongchoice{การเก็บภาษีด้วยสินค้า}
    \end{choices}
  \end{question}
}

\element{general}{
  \begin{question}{Q055}
    Command Economy คืออะไร
    \begin{choices}
        \correctchoice{ระบบเศรษฐกิจที่รัฐบาลควบคุมการผลิตและราคา}
        \wrongchoice{ผู้บริโภคเลือกเองอิสระ}
        \wrongchoice{เอกชนแข่งขันเสรี}
        \wrongchoice{เศรษฐกิจการค้าระหว่างประเทศ}
    \end{choices}
  \end{question}
}

\element{general}{
  \begin{question}{Q056}
    Market Economy คืออะไร
    \begin{choices}
        \wrongchoice{รัฐบาลตัดสินใจการผลิต}
        \correctchoice{ราคาและการผลิตกำหนดโดยอุปสงค์อุปทาน}
        \wrongchoice{อุตสาหกรรมเป็นของรัฐ}
        \wrongchoice{เศรษฐกิจไม่ใช้เงิน}
    \end{choices}
  \end{question}
}

% \element{general}{
%   \begin{question}{Q057}
%     จุดที่มีประสิทธิภาพที่สุดบนเส้น PPF คือจุดใด
%     \begin{choices}
%         \correctchoice{บนเส้นกราฟ (On the frontier)}
%         \wrongchoice{ใต้เส้นกราฟ}
%         \wrongchoice{เหนือเส้นกราฟ}
%         \wrongchoice{จุดใดก็ได้}
%     \end{choices}
%   \end{question}
% }

% \element{general}{
%   \begin{question}{Q058}
%     ประโยชน์ของการค้าขาย (Gain from trading) ตามทฤษฎีคลาสสิกคืออะไร
%     \begin{choices}
%         \correctchoice{ผู้คนสามารถมีสินค้าได้มากกว่าที่ผลิตเอง}
%         \wrongchoice{ผลิตได้เกินเส้น PPF ของตนเอง}
%         \wrongchoice{มีเวลามากขึ้น}
%         \wrongchoice{ไม่มีข้อถูก}
%     \end{choices}
%   \end{question}
% }

% \element{general}{
%   \begin{question}{Q059}
%     "มือที่มองไม่เห็น" (Invisible Hand) หมายถึงอะไร
%     \begin{choices}
%         \wrongchoice{รัฐบาลแทรกแซง}
%         \correctchoice{กลไกตลาดที่ขับเคลื่อนด้วยการแสวงหาผลประโยชน์ส่วนตน}
%         \wrongchoice{ข้อตกลงลับ}
%         \wrongchoice{AI}
%     \end{choices}
%   \end{question}
% }

\element{general}{
  \begin{question}{Q060}
    Collaborative Economy คืออะไร
    \begin{choices}
        \wrongchoice{แลกเปลี่ยนสินค้าด้วยเงิน}
        \wrongchoice{แลกเปลี่ยนสินค้ากับสินค้า}
        \correctchoice{ชื่อเรียกอีกอย่างของ Sharing Economy}
        \wrongchoice{ผิดทุกข้อ}
    \end{choices}
  \end{question}
}

\element{general}{
  \begin{question}{Q061}
    ข้อใด \textbf{ไม่ใช่} ตัวอย่างของแพลตฟอร์ม Collaborative Economy
    \begin{choices}
        \wrongchoice{Lyft}
        \correctchoice{Netflix}
        \wrongchoice{Airbnb}
        \wrongchoice{Uber}
    \end{choices}
  \end{question}
}

\element{general}{
  \begin{question}{Q062}
    Gig Economy คืออะไร
    \begin{choices}
        \correctchoice{ตลาดแรงงานที่มีตำแหน่งงานชั่วคราว พาร์ทไทม์}
        \wrongchoice{เศรษฐกิจฮาร์ดแวร์}
        \wrongchoice{เศรษฐกิจ 1 GB}
        \wrongchoice{ผิดทุกข้อ}
    \end{choices}
  \end{question}
}

\element{general}{
  \begin{question}{Q063}
    ข้อเปรียบเทียบ Sharing Economy vs Gig Economy ข้อใดถูก
    \begin{choices}
        \wrongchoice{ไม่เกี่ยวข้องกัน}
        \wrongchoice{เหมือนกัน}
        \correctchoice{ใช้ทรัพยากรที่ยังใช้ไม่เต็มที่เหมือนกัน แต่ Sharing เน้นสินทรัพย์ ส่วน Gig เน้นบริการ}
        \wrongchoice{ใช้ทรัพยากรที่ยังใช้ไม่เต็มที่เหมือนกัน Sharing เน้นบริการ Gig เน้นสินทรัพย์}
    \end{choices}
  \end{question}
}

\element{general}{
  \begin{question}{Q064}
    Product Complexity มักจะหมายถึงอะไร
    \begin{choices}
        \wrongchoice{ข้อมูลการศึกษา}
        \correctchoice{ข้อมูลที่จำเป็นในการอธิบายเพื่อการทำธุรกรรม}
        \wrongchoice{ปริมาณข้อมูลที่ดึงดูด}
        \wrongchoice{ผิดทุกข้อ}
    \end{choices}
  \end{question}
}

\element{general}{
  \begin{question}{Q065}
    สินค้าใดที่มีความซับซ้อนของผลิตภัณฑ์สูงที่สุด
    \begin{choices}
        \wrongchoice{เสื้อผ้า}
        \wrongchoice{น้ำดื่ม}
        \wrongchoice{บ้าน}
        \correctchoice{รถยนต์}
    \end{choices}
  \end{question}
}

\element{general}{
  \begin{question}{Q066}
    Asset Specificity คืออะไร
    \begin{choices}
        \wrongchoice{ข้อมูลการศึกษา}
        \wrongchoice{ข้อมูลธุรกรรม}
        \correctchoice{ระดับความเฉพาะเจาะจงของสินทรัพย์}
        \wrongchoice{ผิดทุกข้อ}
    \end{choices}
  \end{question}
}

\element{general}{
  \begin{question}{Q067}
    MYB Framework พยายามอธิบายอะไร
    \begin{choices}
        \wrongchoice{กิจกรรมภายในบริษัทขยายตัว}
        \correctchoice{กิจกรรมของตลาดขยายตัว}
        \wrongchoice{มีแพลตฟอร์มมากขึ้น}
        \wrongchoice{ผิดทุกข้อ}
    \end{choices}
  \end{question}
}

\element{general}{
  \begin{question}{Q068}
    สินค้าประเภทใดเหมาะกับ Collaborative Economy
    \begin{choices}
        \wrongchoice{มูลค่าสูง ใช้บ่อย}
        \correctchoice{มูลค่าสูง ใช้น้อย}
        \wrongchoice{มูลค่าต่ำ ใช้บ่อย}
        \wrongchoice{มูลค่าต่ำ ใช้น้อย}
    \end{choices}
  \end{question}
}

\element{general}{
  \begin{question}{Q069}
    สูตรคำนวณมูลค่าการเช่าแฝง (Latent Rental Value) คือ
    \begin{choices}
        \correctchoice{\% Idle time $\times$ Product value}
        \wrongchoice{\% Used time $\times$ Product value}
        \wrongchoice{\% Idle time $\times$ Time value}
        \wrongchoice{ผิดทุกข้อ}
    \end{choices}
  \end{question}
}

\element{general}{
  \begin{question}{Q070}
    ประโยชน์ของ Collaborative Economy คืออะไร
    \begin{choices}
        \correctchoice{การใช้สินทรัพย์ที่ว่างเว้นให้เกิดประโยชน์สูงสุด}
        \wrongchoice{คนเชื่อมต่อกัน}
        \wrongchoice{ทำงานเป็นทีม}
        \wrongchoice{ผิดทุกข้อ}
    \end{choices}
  \end{question}
}

\element{general}{
  \begin{question}{Q071}
    ทำไมเศรษฐกิจฐานความสนใจ (Attention Economy) ถึงสำคัญ
    \begin{choices}
        \wrongchoice{ต้นทุนผลิตสูง}
        \wrongchoice{ขายข้อมูล}
        \correctchoice{เพราะความสนใจมนุษย์เป็นทรัพยากรหายาก}
        \wrongchoice{ถูกทุกข้อ}
    \end{choices}
  \end{question}
}

\element{general}{
  \begin{question}{Q072}
    ข้อดีของการได้ตำแหน่งสูงใน Search Engine คืออะไร
    \begin{choices}
        \wrongchoice{ค่าโฆษณาถูก}
        \wrongchoice{รูปแบบดีกว่า}
        \correctchoice{มีโอกาสที่ลูกค้าจะคลิกและซื้อมากขึ้น}
        \wrongchoice{คนเห็นไม่ว่าจะค้นคำว่าอะไร}
    \end{choices}
  \end{question}
}

\element{general}{
  \begin{question}{Q073}
    วัตถุประสงค์หลักของ \textbf{SEO} (Search Engine Optimization) คืออะไร
    \begin{choices}
        \wrongchoice{ค้นหาเร็วขึ้น}
        \correctchoice{เพื่อให้เว็บไซต์ติดอันดับดี}
        \wrongchoice{ลดค่าโฆษณา}
        \wrongchoice{หาพันธมิตร}
    \end{choices}
  \end{question}
}

\element{general}{
  \begin{question}{Q074}
    ข้อใดเปรียบเทียบความแตกต่างระหว่าง \textbf{PPC} (Pay-Per-Click) และ \textbf{SEO} (Search Engine Optimization) ได้ถูกต้อง
    \begin{choices}
        \wrongchoice{PPC เน้นผลลัพธ์ระยะยาวและยั่งยืนกว่า SEO}
        \wrongchoice{SEO ต้องจ่ายเงินค่าโฆษณาต่อคลิกให้กับ Search Engine}
        \correctchoice{PPC เห็นผลลัพธ์ทันที ส่วน SEO ต้องใช้เวลาเพื่อสร้าง Traffic}
        \wrongchoice{SEO สามารถการันตีอันดับ 1 บน Google ได้แน่นอน}
    \end{choices}
  \end{question}
}

\element{general}{
  \begin{question}{Q075}
    ข้อใด \textbf{ไม่จริง} เกี่ยวกับ Affiliate Marketing
    \begin{choices}
        \wrongchoice{เจ้าของติดตามผลได้}
        \wrongchoice{ให้รางวัลเมื่อมีผลงาน}
        \wrongchoice{Influencer ได้เงินโดยไม่ต้องเป็นเจ้าของสินค้า}
        \correctchoice{เงินที่ได้ไม่ขึ้นกับผลงาน}
    \end{choices}
  \end{question}
}

\element{general}{
  \begin{question}{Q076}
    ตัวอย่างของ Contextual Marketing คือ
    \begin{choices}
        \wrongchoice{ดริฟท์รถ}
        \wrongchoice{หมูเด้ง}
        \correctchoice{ได้รับ SMS คูปองเมื่อเดินเข้าห้าง}
        \wrongchoice{โฆษณาดราม่า}
    \end{choices}
  \end{question}
}

\element{general}{
  \begin{question}{Q077}
    จากตาราง Hawk-Dove: หากสัตว์ตัวที่ 1 เลือก \textbf{Dove} \\
    การตอบโต้ที่ดีที่สุด (Best Response) ของสัตว์ตัวที่ 2 คืออะไร

    \begin{center}
      \footnotesize % ตัวหนังสือเล็ก
      \setlength{\tabcolsep}{3pt} % บีบตาราง
      \begin{tabular}{c|cc}
         & \multicolumn{2}{c}{\textbf{Animal 2}} \\
        \textbf{Animal 1} & \textbf{Hawk} & \textbf{Dove} \\ \hline
        \textbf{Hawk} & (-5, -5) & (10, 0) \\
        \textbf{Dove} & (0, 10) & (5, 5)
      \end{tabular}
    \end{center}

    \begin{choices}
        \wrongchoice{Dove}
        \correctchoice{Hawk}
        \wrongchoice{ทั้งคู่}
        \wrongchoice{ไม่มี}
    \end{choices}
  \end{question}
}

\element{general}{
  \begin{question}{Q078}
    จากตาราง Partner: กลยุทธ์ใดของคุณคือ \textbf{Dominated Strategy} \\
    {\footnotesize (คือทางเลือกที่ได้คะแนนน้อยกว่าเสมอ ไม่ว่าเพื่อนจะเลือกอะไร)}

    \begin{center}
      % ลดขนาดฟอนต์เฉพาะในตาราง
      \footnotesize 
      % บีบระยะห่างระหว่างคอลัมน์ให้ชิดกัน (ปกติ 6pt)
      \setlength{\tabcolsep}{3pt} 
      \begin{tabular}{c|cc}
         & \multicolumn{2}{c}{\textbf{Partner}} \\
        \textbf{You} & \textbf{Presentation} & \textbf{Exam} \\ \hline
        \textbf{Presentation} & (8, 8) & (0, 10) \\
        \textbf{Exam} & (10, 0) & (2, 2)
      \end{tabular}
    \end{center}

    \begin{choices}
        \correctchoice{Presentation}
        \wrongchoice{Exam}
        \wrongchoice{ทั้งคู่}
        \wrongchoice{ไม่มี}
    \end{choices}
  \end{question}
}

\element{general}{
  \begin{question}{Q079}
    พิจารณาตารางผลตอบแทน (Payoff Matrix) ของสงครามการค้า (Trade War) ระหว่างประเทศ A และ B ด้านล่างนี้
    
    \begin{center}
    \begin{tabular}{c|c|c}
       & \textbf{B: Free Trade} & \textbf{B: Tariff} \\ \hline
      \textbf{A: Free Trade} & (10, 10) & (-5, 15) \\ \hline
      \textbf{A: Tariff} & (15, -5) & (0, 0)
    \end{tabular}
    \end{center}
    
    คู่กลยุทธ์ใดคือจุดดุลยภาพของแนช (Nash Equilibrium)

    \begin{choices}
        \wrongchoice{Free Trade, Free Trade}
        \correctchoice{Tariff, Tariff}
        \wrongchoice{Free Trade, Tariff}
        \wrongchoice{Tariff, Free Trade}
    \end{choices}
  \end{question}
}

\element{general}{
  \begin{question}{Q080}
    ข้อดีของ Second-price auction คือ
    \begin{choices}
        \correctchoice{ส่งเสริมให้ยื่นราคาตามมูลค่าจริง}
        \wrongchoice{ง่ายกว่า}
        \wrongchoice{ไม่มีอคติ}
        \wrongchoice{ใช้ได้หลากหลาย}
    \end{choices}
  \end{question}
}

% \element{general}{
%   \begin{question}{Q081}
%     ใครได้ประโยชน์จากกฎการประมูลแบบ Vickrey (Second-price)
%     \begin{choices}
%         \wrongchoice{ผู้จัดประมูล}
%         \correctchoice{ผู้เข้าประมูล}
%         \wrongchoice{คู่แข่ง}
%         \wrongchoice{เท่ากัน}
%     \end{choices}
%   \end{question}
% }

\element{general}{
  \begin{question}{Q082}
    ข้อใด \textbf{ไม่ใช่} ปัจจัยหลักของ Ad Rank ใน Google
    \begin{choices}
        \wrongchoice{เงินประมูล}
        \wrongchoice{คุณภาพของหน้า Landing Page}
        \wrongchoice{รูปแบบโฆษณา}
        \correctchoice{ความมีสีสัน}
    \end{choices}
  \end{question}
}

% \element{general}{
%   \begin{question}{Q083}
%     Trust Economy คืออะไร
%     \begin{choices}
%         \wrongchoice{เศรษฐกิจกฎระเบียบ}
%         \wrongchoice{เศรษฐกิจไร้เงินสด}
%         \correctchoice{เศรษฐกิจที่ความเชื่อใจเป็นสินทรัพย์สำคัญ}
%         \wrongchoice{ถูกทุกข้อ}
%     \end{choices}
%   \end{question}
% }

% \element{general}{
%   \begin{question}{Q084}
%     ทำไมโดรนพาณิชย์ถึงใช้เยอะในสงครามรัสเซีย-ยูเครน
%     \begin{choices}
%         \wrongchoice{ออกแบบเพื่อการทหาร}
%         \correctchoice{ราคาถูก หาง่าย ดัดแปลงง่าย}
%         \wrongchoice{ต้องฝึกฝนสูง}
%         \wrongchoice{บินได้นานกว่า}
%     \end{choices}
%   \end{question}
% }

% \element{general}{
%   \begin{question}{Q085}
%     Reshoring คืออะไร
%     \begin{choices}
%         \wrongchoice{ย้ายไปประเทศเพื่อนบ้าน}
%         \correctchoice{ย้ายฐานผลิตกลับประเทศต้นกำเนิด}
%         \wrongchoice{จ้างผลิตที่อื่น}
%         \wrongchoice{ตั้งโรงงานใหม่ต่างประเทศ}
%     \end{choices}
%   \end{question}
% }

% \element{general}{
%   \begin{question}{Q086}
%     Friend-shoring คืออะไร
%     \begin{choices}
%         \wrongchoice{จ้างแรงงานถูก}
%         \correctchoice{ย้ายฐานผลิตไปประเทศพันธมิตรที่มีเสถียรภาพ}
%         \wrongchoice{ผลิตในประเทศเท่านั้น}
%         \wrongchoice{ลดภาษี}
%     \end{choices}
%   \end{question}
% }

\element{general}{
  \begin{question}{Q087}
    ข้อตกลง Bretton Woods สร้างโครงสร้างพื้นฐานใด
    \begin{choices}
        \wrongchoice{EU}
        \correctchoice{ระบบอัตราแลกเปลี่ยนคงที่และสถาบันการเงิน}
        \wrongchoice{กองทัพ}
        \wrongchoice{UN}
    \end{choices}
  \end{question}
}

\element{general}{
  \begin{question}{Q088}
    สิ่งใดช่วยพยุงอำนาจดอลลาร์สหรัฐหลัง Bretton Woods
    \begin{choices}
        \wrongchoice{Silk road}
        \correctchoice{Petro dollars}
        \wrongchoice{De-Globalization}
        \wrongchoice{Digital Economy}
    \end{choices}
  \end{question}
}

\element{general}{
  \begin{question}{Q089}
    เป้าหมายหนึ่งของ BRICS คือ
    \begin{choices}
        \wrongchoice{สนับสนุนตะวันตก}
        \wrongchoice{พึ่งพาดอลลาร์}
        \correctchoice{สร้างความร่วมมือเศรษฐกิจเกิดใหม่และลดการใช้ดอลลาร์}
        \wrongchoice{พันธมิตรทางทหาร}
    \end{choices}
  \end{question}
}

% \element{general}{
%   \begin{question}{Q090}
%     แรงขับเคลื่อนความขัดแย้งของมหาอำนาจปัจจุบันคือ
%     \begin{choices}
%         \wrongchoice{อุดมการณ์}
%         \wrongchoice{ดินแดน}
%         \correctchoice{การควบคุมทรัพยากรพลังงานและห่วงโซ่อุปทาน AI}
%         \wrongchoice{วัฒนธรรม}
%     \end{choices}
%   \end{question}
% }

\element{general}{
  \begin{question}{Q-GameTheory-Ad}
    พิจารณากรณีบริษัทเครื่องดื่ม 2 ราย (A และ B) ต้องตัดสินใจว่าจะ ``โฆษณา'' หรือ ``ไม่โฆษณา''
    
    \begin{itemize}
        \item หากทั้งคู่ ไม่โฆษณา ได้กำไรรายละ 50 ล้าน (ประหยัดงบ)
        \item หากทั้งคู่ โฆษณา ได้กำไรรายละ 20 ล้าน (หักล้างกันเองและเสียค่าโฆษณา)
        \item หาก ฝ่ายหนึ่งโฆษณา แต่อีกฝ่ายไม่ ผู้โฆษณาได้ 80 ล้าน, ผู้ไม่โฆษณาเหลือ 0
    \end{itemize}
    
    ในสถานการณ์นี้ จุดดุลยภาพของแนช (Nash Equilibrium) คือข้อใด และส่งผลอย่างไร
    
    \begin{choices}
        \correctchoice{ทั้งคู่เลือก โฆษณา}
        \wrongchoice{ทั้งคู่เลือก ไม่โฆษณา}
        \wrongchoice{ฝ่ายหนึ่งโฆษณา อีกฝ่ายไม่โฆษณา}
        \wrongchoice{ไม่มีข้อใดเป็นดุลยภาพที่แน่นอน}
    \end{choices}
  \end{question}
}

% ======================================================================
% GROUP 2: QUIZZ (ข้อสอบใหม่จากไฟล์ PDF)
% ======================================================================

\element{quizz}{
  \begin{question}{QZ001}
    ข้อใดต่อไปนี้ถือเป็นองค์ประกอบพื้นฐานของเศรษฐกิจดิจิทัล
    \begin{choices}
        \wrongchoice{เศรษฐกิจบริการ}
        \wrongchoice{เศรษฐกิจอุตสาหกรรม}
        \wrongchoice{เศรษฐกิจการเกษตร}
        \correctchoice{เศรษฐกิจแพลตฟอร์ม}
    \end{choices}
  \end{question}
}

\element{quizz}{
  \begin{question}{QZ002}
    สูตร ``GDP จากด้านรายจ่าย'' ที่ถูกต้องคือข้อใด
    \begin{choices}
        \wrongchoice{$C+I+G-NX+eK$}
        \wrongchoice{$C+I+G+NX+M$}
        \correctchoice{$C+I+G+NX$}
        \wrongchoice{$C+I+NX$}
    \end{choices}
  \end{question}
}

\element{quizz}{
  \begin{question}{QZ003}
    ตาม Solow model เหตุใดการเพิ่มทุนและแรงงานอย่างต่อเนื่องจึงไม่ยั่งยืน
    \begin{choices}
        \wrongchoice{ผลผลิตจะลดลงในที่สุด}
        \correctchoice{ค่าเสื่อมราคาจะสูงกว่าระดับผลผลิต การลงทุน}
        \wrongchoice{ทุนและแรงงานมีข้อจำกัด}
        \wrongchoice{ต้นทุนส่วนเพิ่มจะเข้าใกล้ศูนย์}
    \end{choices}
  \end{question}
}

\element{quizz}{
  \begin{question}{QZ004}
    Solow model เป็นแบบจำลองที่อธิบายเรื่องใด
    \begin{choices}
        \wrongchoice{ความไม่เท่าเทียมกันของรายได้}
        \wrongchoice{ความแตกต่างระหว่างประเทศ}
        \correctchoice{การเติบโตทางเศรษฐกิจ}
        \wrongchoice{ถูกทุกข้อ}
    \end{choices}
  \end{question}
}

\element{quizz}{
  \begin{question}{QZ005}
    จงเลือกชุดของสัญลักษณ์ที่ตรงกับ Solow model และลำดับต่อไปนี้ human capital, physical capital, ideas
    \begin{choices}
        \wrongchoice{eL, eK, eA}
        \correctchoice{eL, K, A}
        \wrongchoice{eL, eK, A}
        \wrongchoice{eK, L, A}
    \end{choices}
  \end{question}
}

\element{quizz}{
  \begin{question}{QZ006}
    ข้อได้เปรียบสำคัญของระบบดิจิทัลเมื่อเทียบกับระบบอนาล็อกคืออะไร
    \begin{choices}
        \wrongchoice{การกำหนดมาตรฐานช่วยให้สามารถใช้งานร่วมกันได้}
        \wrongchoice{สามารถผลิตอุปกรณ์ในขนาดที่เล็กลงได้}
        \wrongchoice{สามารถส่งสัญญาณได้ทันทีไปยังทุกสถานที่}
        \correctchoice{ถูกทุกข้อ}
    \end{choices}
  \end{question}
}

% \element{quizz}{
%   \begin{question}{QZ007}
%     การแปลงสัญญาณเสียงอนาล็อกเป็นดิจิทัลเรียกว่าอะไร
%     \begin{choices}
%         \wrongchoice{Encoding}
%         \correctchoice{Sampling}
%         \wrongchoice{Rendering}
%         \wrongchoice{Filtering}
%     \end{choices}
%   \end{question}
% }

% \element{quizz}{
%   \begin{question}{QZ008}
%     ความต่างศักย์ไฟฟ้า (Voltage) ที่ส่งสัญญาณสถานะ "เปิด" (Logic 1) ในระบบดิจิทัล (TTL) โดยทั่วไปมีค่าประมาณเท่าใด
%     \begin{choices}
%         \wrongchoice{1 โวลต์}
%         \wrongchoice{2 โวลต์}
%         \wrongchoice{3 โวลต์}
%         \correctchoice{5 โวลต์}
%     \end{choices}
%   \end{question}
% }

\element{quizz}{
  \begin{question}{QZ009}
    ต้องใช้กี่บิตในการแสดงพิกเซลเดียวบนจอภาพขาวดำ (Black \& White)
    \begin{choices}
        \wrongchoice{0}
        \correctchoice{1}
        \wrongchoice{2}
        \wrongchoice{4}
    \end{choices}
  \end{question}
}

\element{quizz}{
  \begin{question}{QZ010}
    ต้องใช้กี่บิตในการแสดงพิกเซลเดียวบนจอภาพ 256 สี
    \begin{choices}
        \wrongchoice{5}
        \wrongchoice{6}
        \wrongchoice{7}
        \correctchoice{8}
    \end{choices}
  \end{question}
}

\element{quizz}{
  \begin{question}{QZ011}
    ``Spatial Data Management System'' (SDMS) เกี่ยวข้องกับอะไร
    \begin{choices}
        \correctchoice{การจัดการข้อมูลเชิงพื้นที่}
        \wrongchoice{การบีบอัดข้อมูลเสียง}
        \wrongchoice{การทำแผนที่สมอง}
        \wrongchoice{การวิเคราะห์ไฟล์รูปภาพ}
    \end{choices}
  \end{question}
}

\element{quizz}{
  \begin{question}{QZ012}
    ข้อใดต่อไปนี้เป็นตัวอย่างของรูปแบบธุรกิจ B2B
    \begin{choices}
        \correctchoice{บริษัทผลิตวัตถุดิบขายให้โรงงานผลิตสินค้า}
        \wrongchoice{ร้านค้าออนไลน์ขายของให้ผู้บริโภค}
        \wrongchoice{ลูกค้าขายของมือสองให้กันเอง}
        \wrongchoice{ผู้บริโภคจ้างนักออกแบบอิสระ}
    \end{choices}
  \end{question}
}

\element{quizz}{
  \begin{question}{QZ013}
    บริษัทเอกชนทำระบบจัดเก็บภาษีออนไลน์ให้กับกรมสรรพากร ถือเป็นธุรกรรมแบบใด
    \begin{choices}
        \correctchoice{B2G}
        \wrongchoice{B2C}
        \wrongchoice{C2G}
        \wrongchoice{C2C}
    \end{choices}
  \end{question}
}

\element{quizz}{
  \begin{question}{QZ014}
    ``ฟรีแลนซ์ขายภาพถ่ายเชิงพาณิชย์ให้บริษัทผ่านแพลตฟอร์ม'' เข้ากับประเภทใด
    \begin{choices}
        \wrongchoice{B2B}
        \correctchoice{C2B}
        \wrongchoice{C2C}
        \wrongchoice{B2C}
    \end{choices}
  \end{question}
}

\element{quizz}{
  \begin{question}{QZ015}
    เมื่อเปรียบเทียบกับโมเดลธุรกิจแบบดั้งเดิม ต้นทุนใดที่สามารถลดลงได้อย่างมากจากการขายตรง
    \begin{choices}
        \wrongchoice{การผลิต}
        \correctchoice{การกระจายสินค้าหรือคนกลาง}
        \wrongchoice{การขายส่ง}
        \wrongchoice{การโฆษณา}
    \end{choices}
  \end{question}
}

\element{quizz}{
  \begin{question}{QZ016}
    ข้อใดต่อไปนี้เป็นตัวอย่างของ G2C
    \begin{choices}
        \wrongchoice{การจัดซื้อจัดจ้างภาครัฐผ่านระบบอิเล็กทรอนิกส์}
        \correctchoice{แอปมือถือให้ประชาชนต่ออายุใบขับขี่และรับใบอนุญาตดิจิทัล}
        \wrongchoice{ระบบแลกเปลี่ยนข้อมูลระหว่างหน่วยงานรัฐ}
        \wrongchoice{แพลตฟอร์มออกหนังสือรับรองส่งออก}
    \end{choices}
  \end{question}
}

\element{quizz}{
  \begin{question}{QZ017}
    ข้อใดต่อไปนี้เป็นแหล่งที่มาของต้นทุนผันแปร (Variable Cost)
    \begin{choices}
        \wrongchoice{ค่าเช่าสำนักงาน}
        \correctchoice{วัตถุดิบ}
        \wrongchoice{การตรวจสอบบัญชี}
        \wrongchoice{ประกันอัคคีภัย}
    \end{choices}
  \end{question}
}

% \element{quizz}{
%   \begin{question}{QZ018}
%     ความแตกต่างหลักระหว่าง ผู้บริโภคใน C2C และ ธุรกิจใน B2C คืออะไร
%     \begin{choices}
%         \correctchoice{C2C มักเป็นผู้ขายชั่วคราว/รายย่อย}
%         \wrongchoice{C2C ขายออนไลน์เท่านั้น}
%         \wrongchoice{C2C ไม่ได้ใช้แพลตฟอร์ม}
%         \wrongchoice{ไม่มีความแตกต่าง}
%     \end{choices}
%   \end{question}
% }

\element{quizz}{
  \begin{question}{QZ019-Concept}
    ในทางเศรษฐศาสตร์ บริการ Streaming เช่น Netflix จัดเป็น Club Goods (สินค้าสโมสร) เนื่องจากเหตุผลใด
    
    \begin{choices}
        \wrongchoice{Non-excludable, Rivalrous (ใครก็ใช้ได้ แต่แย่งกันใช้)}
        \wrongchoice{Excludable, Rivalrous (ต้องจ่ายเงิน และแย่งกันใช้)}
        \correctchoice{Excludable, Non-rivalrous (ต้องจ่ายเงิน แต่ไม่แย่งกันใช้)}
        \wrongchoice{Non-excludable, Non-rivalrous (ใครก็ใช้ได้ และไม่แย่งกันใช้)}
    \end{choices}
  \end{question}
}

\element{quizz}{
  \begin{question}{QZ020}
    Airbnb ให้บริการประเภทใด
    \begin{choices}
        \wrongchoice{ขนส่ง}
        \wrongchoice{วีดีโอ}
        \correctchoice{ที่พักหรือโรงแรม}
        \wrongchoice{การบิน}
    \end{choices}
  \end{question}
}

% \element{quizz}{
%   \begin{question}{QZ021}  <-- DELETED: Duplicate concept of DE001 (Hard) and Q042 (General)
%     ประโยชน์ทางเศรษฐกิจของสินค้าสารสนเทศ (Information Goods) บนแพลตฟอร์มคืออะไร
%     \begin{choices}
%         \wrongchoice{รายได้เพิ่มขึ้น}
%         \correctchoice{ต้นทุนส่วนเพิ่มต่ำ}
%         \wrongchoice{ลูกค้าเพิ่มขึ้น}
%         \wrongchoice{ไม่จำเป็นต้องโฆษณา}
%     \end{choices}
%   \end{question}
% }

\element{quizz}{
  \begin{question}{QZ022}
    จุดใดที่ ต้นทุนการผลิตเฉลี่ย (Average Cost) มักจะต่ำที่สุด
    \begin{choices}
        \correctchoice{ที่จุดตัดกันระหว่าง Average Total Cost (ATC) และ Marginal Cost (MC)}
        \wrongchoice{เมื่อ Marginal Cost = 0}
        \wrongchoice{เมื่อความชันของ Marginal Cost = 0}
        \wrongchoice{เมื่อความชันของ Marginal Cost = infinity}
    \end{choices}
  \end{question}
}

\element{quizz}{
  \begin{question}{QZ023}
    ต้นทุนที่สำคัญที่สุดใน Cross-border e-commerce (CBEC) โดยทั่วไปคืออะไร
    \begin{choices}
        \wrongchoice{ค่าโฆษณา}
        \wrongchoice{ค่าพัฒนาแพลตฟอร์มมือถือ}
        \wrongchoice{ต้นทุนจม}
        \correctchoice{ต้นทุนด้านโลจิสติกส์}
    \end{choices}
  \end{question}
}

\element{quizz}{
  \begin{question}{QZ024}
    CAPEX (Capital Expense) หมายถึงค่าใช้จ่ายประเภทใด
    \begin{choices}
        \wrongchoice{ค่าเชื้อเพลิงรายเดือน}
        \wrongchoice{ค่าเดินเครื่องตามปกติ}
        \correctchoice{เงินลงทุนก้อนใหญ่สำหรับสร้างหรือติดตั้งโรงไฟฟ้า}
        \wrongchoice{ค่าปรับสิ่งแวดล้อม}
    \end{choices}
  \end{question}
}

\element{quizz}{
  \begin{question}{QZ025}
    OPEX (Operating Expense) คือค่าใช้จ่ายใด
    \begin{choices}
        \wrongchoice{ภาษีที่ดิน}
        \correctchoice{ค่าเชื้อเพลิง ค่าเดินเครื่อง}
        \wrongchoice{ค่าเสื่อมราคา}
        \wrongchoice{ค่าธรรมเนียมใบอนุญาต}
    \end{choices}
  \end{question}
}

\element{quizz}{
  \begin{question}{QZ026}
    ถ้าต้องการ ลด Greenhouse Gas ภาคไฟฟ้าควรทำอะไร
    \begin{choices}
        \wrongchoice{เพิ่มสัดส่วนถ่านหิน}
        \correctchoice{เพิ่มสัดส่วนพลังงานหมุนเวียน}
        \wrongchoice{ลดกำลังผลิตทั้งหมด}
        \wrongchoice{เพิ่มค่าบำรุงรักษาอย่างเดียว}
    \end{choices}
  \end{question}
}

\element{quizz}{
  \begin{question}{QZ027}
    สมมติว่ามีคนมาจ้างให้คุณเขียนโปรแกรมเป็นโปรเจคสั้น ๆ คุณกำลังอยู่ในระบบเศรษฐกิจแบบใด
    \begin{choices}
        \wrongchoice{Solow Economy}
        \correctchoice{Gig Economy}
        \wrongchoice{Sharing Economy}
        \wrongchoice{ถูกทุกข้อ}
    \end{choices}
  \end{question}
}

\element{quizz}{
  \begin{question}{QZ028}
    สมมติว่าคุณนำที่ว่างหน้าบ้านมาทำที่จอดรถให้คนอื่นเช่า คุณกำลังอยู่ในเศรษฐกิจแบบใด
    \begin{choices}
        \wrongchoice{Solow Economy}
        \wrongchoice{Gig Economy}
        \correctchoice{Sharing Economy}
        \wrongchoice{ถูกทุกข้อ}
    \end{choices}
  \end{question}
}

\element{quizz}{
  \begin{question}{QZ029}
    จีนและไต้หวันขัดแย้งกันด้วยสาเหตุหลักใด
    \begin{choices}
        \wrongchoice{ภาษาไม่เหมือนกัน}
        \correctchoice{นโยบายจีนเดียว (One China Policy)}
        \wrongchoice{ไต้หวันไม่ยอมขายสินค้าให้}
        \wrongchoice{ไม่มีข้อใดถูก}
    \end{choices}
  \end{question}
}

\element{quizz}{
  \begin{question}{QZ030}
    ข้อใดถือเป็น Platform
    \begin{choices}
        \wrongchoice{Web 2.0}
        \correctchoice{Facebook}
        \wrongchoice{Microsoft Word}
        \wrongchoice{Keyboard}
    \end{choices}
  \end{question}
}

\element{quizz}{
  \begin{question}{QZ031}
    ในการประมูลแบบ First-Price Sealed-Bid (FPSBA) ถ้าคุณประเมินราคาจริง 10,000 บาท คุณควรยื่นเท่าไร
    \begin{choices}
        \wrongchoice{มากกว่า 10,000}
        \correctchoice{น้อยกว่า 10,000}
        \wrongchoice{สองเท่าของราคาประเมิน}
        \wrongchoice{10,000}
    \end{choices}
  \end{question}
}

% \element{quizz}{
%   \begin{question}{QZ032} <-- DELETED: Redundant with Q080 (General) which covers the principle of truthful bidding in Second-price auctions
%     ในการประมูลแบบ Second-Price Sealed-Bid (SPSBA) ถ้าคุณประเมินราคาจริง 10,000 บาท คุณควรยื่นเท่าไร
%     \begin{choices}
%         \wrongchoice{มากกว่า 10,000}
%         \wrongchoice{น้อยกว่า 10,000}
%         \wrongchoice{สองเท่าของราคาประเมิน}
%         \correctchoice{10,000}
%     \end{choices}
%   \end{question}
% }

\element{quizz}{
  \begin{question}{QZ033}
    Nash Equilibrium ในบริบทของการประมูลคืออะไร
    \begin{choices}
        \wrongchoice{รอดูผลคู่แข่ง}
        \wrongchoice{คุยกันเพื่อฮั้วราคา}
        \correctchoice{เสนอราคาตามกลยุทธ์ที่ทำให้ตัวเองได้ประโยชน์ที่สุดโดยคำนึงถึงคู่แข่ง}
        \wrongchoice{เสนอราคามากที่สุด}
    \end{choices}
  \end{question}
}

\element{quizz}{
  \begin{question}{QZ034-Business}
    เหตุใดแพลตฟอร์มโซเชียลมีเดียจึงปรับอัลกอริทึมให้ลดการมองเห็นของโพสต์ขายสินค้าทั่วไปลง
    \begin{choices}
        \correctchoice{เพื่อกระตุ้นให้ผู้ขายซื้อโฆษณา}
        \wrongchoice{เพื่อลดภาระของเซิร์ฟเวอร์เก็บข้อมูล}
        \wrongchoice{เพื่อสนับสนุนให้ลูกค้าไปซื้อที่หน้าร้านจริงแทน}
        \wrongchoice{เพื่อจำกัดจำนวนผู้ขายรายใหม่}
    \end{choices}
  \end{question}
}

\element{quizz}{
  \begin{question}{QZ035}
    ระดับการอธิบาย ``สาเหตุสงคราม'' ตามทฤษฎีความสัมพันธ์ระหว่างประเทศ \textbf{ไม่รวม} ข้อใด
    \begin{choices}
        \wrongchoice{Individual level}
        \wrongchoice{State level}
        \wrongchoice{System level}
        \correctchoice{Cultural level}
    \end{choices}
  \end{question}
}

\element{quizz}{
  \begin{question}{QZ036}
    ข้อใดตรงกับ Individual level ในสาเหตุของสงคราม
    \begin{choices}
        \wrongchoice{อำนาจอนาธิปไตยระหว่างประเทศ}
        \correctchoice{ผู้นำตัดสินใจและมีความทะเยอทะยานหรือนิสัยส่วนตัว}
        \wrongchoice{ดุลอำนาจระหว่างรัฐ}
        \wrongchoice{ประเภทระบอบการปกครองในประเทศ}
    \end{choices}
  \end{question}
}

\element{quizz}{
  \begin{question}{QZ037}
    ข้อใดตรงกับแนวคิด Diversionary Theory
    \begin{choices}
        \correctchoice{ผู้นำจุดชนวนความขัดแย้งภายนอกเพื่อกลบปัญหาภายในหรือสร้างเอกภาพในชาติ}
        \wrongchoice{เศรษฐกิจพึ่งพากันมากทำให้สงครามหมดไป}
        \wrongchoice{ประเทศเล็กถ่วงดุลประเทศใหญ่}
        \wrongchoice{ผู้นำก้าวร้าวเพราะนิสัยส่วนตัว}
    \end{choices}
  \end{question}
}

\element{quizz}{
  \begin{question}{QZ038}
    ``Hybrid War'' ในสงครามรัสเซีย-ยูเครนหมายถึงอะไรเป็นหลัก
    \begin{choices}
        \wrongchoice{การรบทางเศรษฐกิจ}
        \correctchoice{ผสมผสานไซเบอร์ ข้อมูลข่าวสาร อิเล็กทรอนิกส์ และโดรน เข้ากับการรบจริง}
        \wrongchoice{การรบทางบก}
        \wrongchoice{การคว่ำบาตรทางการเงิน}
    \end{choices}
  \end{question}
}

% \element{quizz}{
%   \begin{question}{QZ039}
%     ข้อใดอธิบายถูกต้องเกี่ยวกับ Offshoring และ Outsourcing
%     \begin{choices}
%         \wrongchoice{Offshoring = จ้างคนนอกทำงาน}
%         \wrongchoice{Outsourcing = ย้ายงานไปทำต่างประเทศ}
%         \correctchoice{Offshoring = ย้ายฐานการผลิตไปต่างประเทศ / Outsourcing = จ้างคนนอกทำงาน}
%         \wrongchoice{ทั้ง A และ B}
%     \end{choices}
%   \end{question}
% }

\element{quizz}{
  \begin{question}{QZ040}
    กลยุทธ์ ``Dedollarization'' ของ BRICS \textbf{ไม่รวม} ข้อใด
    \begin{choices}
        \wrongchoice{ใช้เงินท้องถิ่นใน NDB}
        \wrongchoice{ระบบชำระเงินข้ามแดนทางเลือก}
        \wrongchoice{สร้างสกุลเงินร่วม}
        \correctchoice{พึ่งพา SWIFT มากขึ้น}
    \end{choices}
  \end{question}
}

% ======================================================================
% GROUP 3: HARD
% ======================================================================


\element{hard}{
  \begin{question}{H005}
    พิจารณาสถานการณ์ Prisoner's Dilemma ในการตั้งราคาสินค้าของบริษัทคู่แข่ง 2 รายดังนี้
    หากทั้งคู่ ``คงราคาสูง'' จะได้กำไรคนละ 10 ล้าน
    หากฝ่ายหนึ่ง ``ลดราคา'' แต่อีกฝ่ายไม่ลด คนลดราคาจะได้กำไร 15 ล้าน (อีกคนเหลือ 0)
    หากทั้งคู่ ``ลดราคา'' จะได้กำไรคนละ 5 ล้าน
    จุดดุลยภาพของแนช (Nash Equilibrium) คือข้อใด และเพราะเหตุใด
    \begin{choices}
        \wrongchoice{คงราคาสูงทั้งคู่ เพราะได้ผลตอบแทนรวมกันมากที่สุด (20 ล้าน)}
        \wrongchoice{ผลัดกันลดราคาคนละปี เพื่อแบ่งปันกำไรสูงสุด}
        \correctchoice{ลดราคาทั้งคู่ เพราะการลดราคาเป็นกลยุทธ์ที่ดีที่สุดเสมอไม่ว่าคู่แข่งจะเลือกอะไร}
        \wrongchoice{ไม่มีจุดดุลยภาพ เพราะต่างฝ่ายต่างไม่ไว้ใจกัน}
    \end{choices}
  \end{question}
}

\element{hard}{
  \begin{question}{H004}
    เหตุใดธุรกิจแพลตฟอร์มดิจิทัล (เช่น Software หรือ Streaming) จึงมักมีโครงสร้างตลาดแบบ ``Winner-takes-all'' หรือมีผู้ผูกขาดรายใหญ่เพียงไม่กี่ราย
    \begin{choices}
        \wrongchoice{เพราะมีต้นทุนหน่วยสุดท้าย (Marginal Cost) ที่สูงมากในการขยายฐานลูกค้า}
        \wrongchoice{เพราะรัฐบาลมักให้สัมปทานแก่บริษัทเทคโนโลยีเพียงรายเดียว}
        \correctchoice{เพราะผลของ Network Effect และต้นทุนหน่วยสุดท้ายที่เข้าใกล้ศูนย์ (Zero Marginal Cost)}
        \wrongchoice{เพราะสินค้าดิจิทัลมีความเสื่อมสภาพเร็ว ทำให้คู่แข่งตามไม่ทัน}
    \end{choices}
  \end{question}
}

\element{hard}{
  \begin{question}{H003}
    สมมติการประมูลโฆษณาแบบ Second-price Auction มีผู้ลงโฆษณา 3 ราย ได้แก่ A B และ C
    ผู้ลงโฆษณา A ประเมินมูลค่าคลิกไว้ที่ 100 บาท B ประเมินไว้ 80 บาท และ C ประเมินไว้ 60 บาท
    หากทุกคนยื่นราคาประมูลตามมูลค่าจริงที่ตนประเมินไว้ (Truthful Bidding) ผลลัพธ์จะเป็นอย่างไร
    \begin{choices}
        \wrongchoice{A ชนะประมูล และจ่าย 100 บาท}
        \wrongchoice{A ชนะประมูล และจ่าย 60 บาท}
        \correctchoice{A ชนะประมูล และจ่าย 80 บาท}
        \wrongchoice{B ชนะประมูล เพราะเสนอราคากลาง ๆ}
    \end{choices}
  \end{question}
}

\element{hard}{
  \begin{question}{DE001}
    ลักษณะโครงสร้างต้นทุน (Cost Structure) ของสินค้าดิจิทัล (เช่น ซอฟต์แวร์ ไฟล์เพลง) แตกต่างจากสินค้าอุตสาหกรรมทั่วไปอย่างไร
    \begin{choices}
        \wrongchoice{ต้นทุนคงที่ (Fixed Cost) ต่ำ ต้นทุนหน่วยสุดท้าย (Marginal Cost) สูง}
        \correctchoice{ต้นทุนคงที่ (Fixed Cost) สูงมาก, ต้นทุนหน่วยสุดท้าย (Marginal Cost) เข้าใกล้ศูนย์}
        \wrongchoice{ต้นทุนเฉลี่ย (Average Cost) เพิ่มขึ้นเมื่อผลิตมากขึ้น}
        \wrongchoice{ไม่มีต้นทุนคงที่ มีแต่ต้นทุนผันแปร}
    \end{choices}
  \end{question}
}